\documentclass[12pt]{article}
\usepackage[shortlabels]{enumitem}
\usepackage{float}
\usepackage{xcolor}
\usepackage[a4paper, top=1.5cm, bottom=1.5cm, left=2cm, right=2cm]{geometry}
\usepackage{parskip}
\usepackage{setspace}
\usepackage{lmodern}
\usepackage[T1]{fontenc}
\usepackage[brazil]{babel}
\usepackage{hyperref}
\usepackage{graphicx}
\usepackage{amsmath}
\usepackage{amssymb}
\usepackage{amsfonts}
\usepackage{dsfont}
\usepackage{float}
\title{Lista 6 - Estatística para Administração - 2025.2 - Gabarito}
\begin{document}
\date{}
\maketitle

1) Uma variável aleatória discreta assume valores em um conjunto contável (enumerável), podendo ser finito ou infinito. Já uma variável aleatória 
contínua assume valores em um conjunto não contável (não enumerável) e necessariamente infinito. Um exemplo de variável aleatória discreta é $X$: Número de pessoas 
que irá votar em um candidato específico nas próximas eleições. Um exemplo de variável contínua é $Y:$ Tempo que um ônibus leva até chegar ao seu 
destino final.  


\vspace{3px}

2) $$\binom{30}{4} = \dfrac{30!}{26!4!} = 27405$$

\vspace{3px}

3) 

    \begin{enumerate}[a)]
        \item Binomial(18; 0,25). Podemos olhar por duas perspectivas: A primeira é que cada pessoa acessa o site e clica ou não clica no botão, podendo 
        ser atribuída uma variável aleatória de Bernoulli a essa situação. Como temos 18 pessoas, o total de visitantes que clicam nada mais é 
        que a soma dessas variáveis. Como é razoável assumir que cada cliente toma sua ação sem conhecimento da ação dos demais, podemos 
        assumir independência e, portanto, representando a soma por $X$, temos que $X \sim Binomial(18; 0,25)$. A segunda e mais simples é verificar que 
        analisar o número de clientes que clicam dentre todos os 18 clientes que acessam o site é exatamente a definição da distribuição Binomial, que 
        modela o número de "sucessos" dentre $n$ tentativas. 
        \item $\binom{18}{5} 0,25^5 \cdot 0,75^{13} = 0,1987$
        \item $\binom{18}{0} 0,25^0 \cdot 0,75^{18} = 0,0056$
        \item $\mathds{P}(X \geq 2) = 1 -\mathds{P}(X<2) =1 - ( \mathds{P}(X=0) + \mathds{P}(X=1)) =  1- 0,0056 - 0,0338 = 0,9606$
    \end{enumerate}

\vspace{3px}

4) $Y \sim Binomial(4; 0,3)$

\vspace{3px}

5) Seja $Z$ uma variável aleatória tal que $Z\sim N(0,1)$, calcule:

\begin{enumerate}[a)]
    \item $0,1587$
    \item $0,6826$
    \item $0,14$
    \item $0,9495$
    \item $0,9948$
    \item $0,9896$
    \item $0$
    \item $0,0215$
    \item $0.0228$
    \item $1$
    \item $0$
\end{enumerate}

\vspace{3px}

6) 

Definamos $Y$: Altura dos homens de 25 anos de idade. 
0.0228
$$\mathds{P}(Y > 188) = \mathds{P}\left(\dfrac{Y - 180}{4} > \dfrac{188 - 180}{4}\right) = \mathds{P}(Z > 2) = 0,0228$$

\vspace{3px}


7) Se X é uma variável aleatória com distribuição Normal com parâmetros $\mu=10$ e $\sigma^2=36$, calcule:

\begin{enumerate}[a)]
    \item $\mathds{P}(X>5) = \mathds{P}\left(\dfrac{X - 10}{6}>\dfrac{5 - 10}{6}\right) = \mathds{P}(Z > -0,83) = 0,7967 $
    \item $\mathds{P}(4<X<16) =  \mathds{P}\left(\dfrac{4 - 10}{6}<\dfrac{X - 10}{6}<\dfrac{16 - 10}{6}\right) = \mathds{P}(-1< Z < 1) = 0,6826 $
    \item $\mathds{P}(X<8) = \mathds{P}\left(\dfrac{X - 10}{6}<\dfrac{8 - 10}{6}\right) = \mathds{P}(Z < -0,33) = 0,3707 $
    \item $\mathds{P}(X<20) = \mathds{P}\left(\dfrac{X - 10}{6}<\dfrac{20 - 10}{6}\right) = \mathds{P}(Z < 1,67) = 0,9525$
    \item $\mathds{P}(X>16)=\mathds{P}\left(\dfrac{X - 10}{6}>\dfrac{16 - 10}{6}\right) = \mathds{P}(Z > 1) = 0,1587$
    \item $\mathds{P}(X>5 | X> 2) = \dfrac{\mathds{P}(X>5) }{\mathds{P}(X>2) } = \dfrac{0,7967}{0,9082} = 0,8772 $
    \item $5,98$
    \item $\mathds{P}(5\leq X \leq 15) = \mathds{P}\left(\dfrac{5 - 10}{6}<\dfrac{X - 10}{6}<\dfrac{15 - 10}{6}\right) = \mathds{P}(-0,83< Z < 0,83) = 0,5934$
\end{enumerate}

\vspace{3px}



8) 
\begin{enumerate}[a)]
    \item $\left[58 - 1,645 \cdot \dfrac{12}{\sqrt{64}}, 58 + 1,645 \cdot \dfrac{12}{\sqrt{64}}\right] = [55,53; 60,46]$
    
    \textcolor{red}{Comentário do Professor}: Aqui vocês poderiam ter usado o valor $1,64$ ou $1,65$ que também estaria certo. 

    Interpretação: Podemos afirmar, com $90\%$ de confiança, que o intervalo obtido contém a média populacional da insatisfação dos colaboradores. 
    De forma geral, podemos afirmar que se fosse realizada a coleta de múltiplas amostras e construído um intervalo de confiança para cada uma delas, 
    aproximadamente 90\% desses intervalos iriam conter a verdadeira média e 10\% não. 
    \item $\left[58 - 1,96 \cdot \dfrac{12}{\sqrt{64}}, 58 + 1,96 \cdot \dfrac{12}{\sqrt{64}}\right] = [55,06 ; 60,94]$
    
    Interpretação: Podemos afirmar, com $95\%$ de confiança, que o intervalo obtido contém a média populacional da insatisfação dos colaboradores. 
    De forma geral, podemos afirmar que se fosse realizada a coleta de múltiplas amostras e construído um intervalo de confiança para cada uma delas, 
    aproximadamente 95\% desses intervalos iriam conter a verdadeira média e 5\% não. 
    \item Aumentar o tamanho amostral implica em aumentar o denominador do fator da margem de erro, reduzindo a margem de erro, o que implica em intervalos 
    com amplitude (tamanho) menores. Já o aumento do nível de confiança implica no aumento do valor de $z_{\alpha/2}$ e esse elemento aparece multiplicando 
    a margem de erro, o que faz com que a amplitude aumente. 
\end{enumerate}

\vspace{3px}

9) 

Se $X\sim Binomial(n,p)$ podemos escrever $X = \sum_{i=1}^n W_i$, onde $W_i, i \in \{1,2, \dots, n\}$ são variáveis aleatórias independentes e todas 
seguindo a distribuição Bernoulli com parâmetro $p$. Podemos dizer isso, pois vimos que uma variável aleatória que segue a distribuição Binomial 
nada mais é do que a soma de variáveis aleatórias Bernoulli Independentes e Igualmente Distribuidas (IID). 

Aplicando as propriedades vistas em sala, temos que:

\begin{itemize}
    \item $\mathds{E}[X] = \mathds{E}\left[\sum_{i=1}^n W_i\right] = \sum_{i=1}^n \mathds{E}[W_i]  = \sum_{i=1}^n p = np$
    \item $Var(X) = Var\left(\sum_{i=1}^n W_i\right) = \sum_{i=1}^n Var(W_i)  = \sum_{i=1}^n p(1-p) = np(1-p)$
\end{itemize}

Note que no ponto 1 não precisamos usar o fato delas serem independentes, enquanto no ponto 2 sim. 
\end{document}